% Fichier complété par tools/complete_missing_corrections.py.
% Source : TEC/TEC-05/14_Sympact/14_Sympact.tex
% Statut : rédigé (5 question(s)).
\begin{corrigebox}[Corrigé rédigé]
\CorrigeQuestion{1}
Le graphe d'analyse contient le bâti, tous les solides
                mobiles et les liaisons du schéma. On y ajoute les poids, l'action motrice,
                la charge résistante et les éventuels ressorts/frottements. Les actions de
                liaison internes sont indiquées par des arêtes entre solides.
\CorrigeQuestion{2}
Les liaisons parfaites entre solides appartenant à
                l'ensemble ont une puissance totale nulle. Seuls un frottement, un ressort,
                un amortisseur ou un actionneur interne apportent une puissance non nulle,
                calculée par le comoment des torseurs d'action et cinématique.
\CorrigeQuestion{3}
On somme les comoments des actions extérieures :
                couples moteurs $C_m\dot q$, forces appliquées $\vec F\cdot\vec V(P)$,
                poids $m\vec g\cdot\vec V_G$ et résistances avec un signe négatif. Les
                réactions des liaisons parfaites avec le bâti ont une puissance nulle au
                point de liaison.
\CorrigeQuestion{4}
L'énergie cinétique de l'ensemble est
                \[E_c=\sum_i\left[\frac12m_iV_{G_i}^2+
                \frac12\vec\Omega_i\cdot\mathbf I_{G_i}\vec\Omega_i\right].\]
                En utilisant la loi entrée--sortie, toutes les vitesses sont exprimées en
                fonction de la vitesse généralisée d'entrée, ce qui conduit à
                $E_c=\tfrac12J_{eq}(q)\dot q^2$.
\CorrigeQuestion{5}
L'énergie cinétique de l'ensemble est
                \[E_c=\sum_i\left[\frac12m_iV_{G_i}^2+
                \frac12\vec\Omega_i\cdot\mathbf I_{G_i}\vec\Omega_i\right].\]
                En utilisant la loi entrée--sortie, toutes les vitesses sont exprimées en
                fonction de la vitesse généralisée d'entrée, ce qui conduit à
                $E_c=\tfrac12J_{eq}(q)\dot q^2$.
\end{corrigebox}
